\documentclass{standalone}
\usepackage{aided_preamble}
\begin{document}

Given \eqref{eq:Pu} we can define the dynamic electron density as

\begin{equation}
\label{eq:dynTrho}
\langle \rho(\mathbf{r}) \rangle_T = \sum_a \int_{-\infty}^{\infty} \rho_a(\mathbf{r}_a - \mathbf{u}_a) P_a(\mathbf{u}) d\mathbf{u}
\end{equation}

where

\begin{equation}
\label{eq:Pa}
    P_j(\mathbf{u}) = (2\pi)^{-3/2} |\mathbf{U}_j|^{-1/2} \exp\left(-\frac{1}{2} \mathbf{u} \mathbf{U}_j^{-1} \mathbf{u}^T\right)
\end{equation}

where $\mathbf{u}_j$ is the displacement of the $j^{th}$ nucleus from its equilibrium position, and $\mathbf{U}_j$ is the
symmetric 3x3 tensor which represents the variance-covariance of the nuclear PDF. The Anisotropic Displacement Parameters (ADP)
which make up the tensor $\mathbf{U}_j$ represent the likelihood for the $j^{th}$ nucleus to be displaced in given directions.
These symmetric 3X3 matrics are the variance-covariance in $\mathbf{u}_j = (x_j, y_j, z_j)$ for nucleus $j$.

The block diagonal elements of the Carteisna MSDA represent the ADPs for each nucleus and can be used to define probability
ellipsoids.

\subsection{equations}

% ---------- convenience macros (no amssymb) ----------
\newcommand{\RR}{\mathbf R}          % boldface R instead of \mathbb{R}
\newcommand{\rvec}{\mathbf r}
\newcommand{\svec}{\mathbf s}
\newcommand{\A}{\mathbf A}
\newcommand{\B}{\mathbf B}
\newcommand{\Cvec}{\mathbf C}
\newcommand{\Uii}{\mathbf U_{ii}}
\newcommand{\Uij}{\mathbf U_{ij}}
\newcommand{\Sig}{\boldsymbol\Sigma_{ij}}
\newcommand{\Lam}{\boldsymbol\Lambda_{ij}}
\newcommand{\alphaa}{\alpha}         % = 2·ex(i)
\newcommand{\betab}{\beta}           % = 2·ex(j)
\newcommand{\gammaab}{\gamma}        % = α+β
\newcommand{\poly}{\mathcal P_{\mu\nu}}

\section*{Atom–resolved dynamic electron density}

\paragraph{Definition.}
\[
\rho_{\mathrm{dyn},i}(\rvec)=
\int_{\RR^{3}}
\rho_{\mathrm{stat},i}(\rvec-\svec)\;
p_i(\svec)\,d^{3}\!s,
\qquad
p_i(\svec)=
\frac{\exp\!\bigl[-\tfrac12\,\svec^{\mathsf T}\Uii^{-1}\svec\bigr]}
     {(2\pi)^{3/2}\sqrt{\det\Uii}} .
\tag{1}
\]

\paragraph{Static atomic density in a contracted Gaussian basis.}
\[
\rho_{\mathrm{stat},i}(\rvec)=
\sum_{\substack{\mu\in i \\ \nu}}
P_{\mu\nu}\;
g_\mu(\rvec)\,g_\nu(\rvec),
\quad
g_\mu(\rvec)=
N_\mu\,
(\rvec-\A)^{\mathbf n_\mu}\,
e^{-\alpha_\mu\lVert\rvec-\A\rVert^{2}} .
\tag{2}
\]

\paragraph{Effective covariance for one primitive pair.}
With doubled exponents
\(
\alphaa = 2\alpha_\mu,\;
\betab  = 2\alpha_\nu,\;
\gammaab=\alphaa+\betab
\)
and centres \(\A,\B\), the harmonic convolution yields
\[
\boxed{
\Sig=
\frac{\mathbf I}{2\gammaab}
+\frac{\alphaa^{2}\mathbf U_{ii}
       +\betab^{2}\mathbf U_{jj}
       +\alphaa\betab(\Uij+\Uij^{\mathsf T})}
      {\gammaab^{2}},
\quad
\Lam=\Sig^{-1}}
\tag{3}
\]

\paragraph{Thermally smeared primitive product.}
Let
\(
\displaystyle
\Cvec=\frac{\alphaa\A+\betab\B}{\gammaab}
\)

Then for $s$-type primitives
\[
[g_\mu g_\nu]_{\text{dyn}}(\rvec)=
N_{\mu\nu}\,
\exp\!\Bigl[-\tfrac{\alphaa\betab}{\gammaab}\lVert\A-\B\rVert^{2}\Bigr]
\exp\!\Bigl[-\tfrac12(\rvec-\Cvec)^{\!\mathsf T}\Lam(\rvec-\Cvec)\Bigr],
\tag{4}
\]
where \(N_{\mu\nu}=N_\mu N_\nu\,(\det\Sig)^{-1/2}\).

\paragraph{Including Cartesian angular momenta.}

TODO: Show how to include angular momenta in the dynamic density.

\end{document}
